\scanpage{87}
% archive foliation: 86

Considérons une courbe elliptique relative $E/S$, alors les \deleted{données}\inserted{épinglages} de Legendre $\kappa:\deleted{\unsure{T}}\Delta_{E/S}\simeq\underline{O}_{\Sigma\deleted{\unsure{J}}}(1)_t^{\unsure{\sharp}}$ forment les sections d'un torseur sous $\unsure{\boldsymbol{\mu}|\mu}_{4S}$ --- le \uline{torseur\inserted{$\mathcal{L}_{E/S}$} des épinglages de Legendre de} $E/S$. (Structure de torseur par multiplication scalaire de $\unsure{\boldsymbol{\mu}|\mu}_4\subset G_m\ldots$).

On \deleted{peut le}\inserted{peut} décrire ce torseur aussi de la façon suivante, sans faire intervenir explicitement $\underline{O}_{\Sigma}(1)_t^{\unsure{\sharp}}$, mais seulement $\underline{O}_{\Sigma}(1)_t$. Considérons les isom
\[
(40)\qquad \kappa':\Delta_{E/S}\simeq\underline{O}_{\Sigma}(1)_t
\]
tels que $\kappa'^{\otimes2}$, interprété comme
\[
(41)\qquad \deleted{\unsure{\kappa'^{\otimes2}}}T_{\Sigma\deleted{\unsure{J}},t}\simeq T_{\Sigma,t}\wedge_{\{\pm1\}}\deleted{\unsure{\unsure{\boldsymbol{\mu}|\mu}_4^*}}\underline{\omega}_{\underline{J}}\deleted{\unsure{\}}}
\]
provienne d'une section de $\underline{\omega}_{\underline{J}}$\inserted{(évidemment unique)}. Ces isom. forment \deleted{un torseur}\inserted{les sections} d'un torseur sous $\unsure{\boldsymbol{\mu}|\mu}_{4S}$, soit $Q_{E/S}$ (qu'on peut interpréter comme le schéma des sommets d'un
